% ═══════════════════════════════════════════════════════════════
% ARBEITSBLATT: Pretérito Indefinido (Einführung -ar Verben)
% ═══════════════════════════════════════════════════════════════
% Sequenz: 08 - Las vacaciones
% Stunde: 08c - Pretérito Indefinido Einführung
% ═══════════════════════════════════════════════════════════════

\documentclass[11pt, a4paper]{article}

% ═══════════════════════════════════════════════════════════════
% SW-MODUS TOGGLE
% ═══════════════════════════════════════════════════════════════
\newif\ifswmodus
\swmodusfalse

% ═══════════════════════════════════════════════════════════════
% PAKETE
% ═══════════════════════════════════════════════════════════════

\usepackage[utf8]{inputenc}
\usepackage[T1]{fontenc}
\usepackage[ngerman]{babel}
\usepackage{helvet}
\renewcommand{\familydefault}{\sfdefault}

\usepackage[margin=2cm]{geometry}
\usepackage{enumitem}
\usepackage{tabularx}
\usepackage{booktabs}
\usepackage{multirow}

\usepackage{xcolor}
\usepackage{framed}
\usepackage{tcolorbox}
\tcbuselibrary{skins}

\usepackage{fancyhdr}
\usepackage{lastpage}
\usepackage{needspace}

% ═══════════════════════════════════════════════════════════════
% FARBEN
% ═══════════════════════════════════════════════════════════════

\definecolor{spanisch-color}{HTML}{c41e3a}
\definecolor{grau-color}{HTML}{888888}
\definecolor{hellgrau-color}{HTML}{f5f5f5}
\definecolor{orange-color}{HTML}{b35c00}

\definecolor{sw-dark}{HTML}{000000}
\definecolor{sw-gray}{HTML}{666666}
\definecolor{sw-white}{HTML}{ffffff}

\ifswmodus
    \colorlet{spanisch}{sw-dark}
    \colorlet{grau}{sw-gray}
    \colorlet{hellgrau}{sw-white}
    \colorlet{orange}{sw-dark}
\else
    \colorlet{spanisch}{spanisch-color}
    \colorlet{grau}{grau-color}
    \colorlet{hellgrau}{hellgrau-color}
    \colorlet{orange}{orange-color}
\fi

\newcommand{\fachfarbe}{spanisch}

% ═══════════════════════════════════════════════════════════════
% VARIABLEN
% ═══════════════════════════════════════════════════════════════

\newcommand{\thema}{Pretérito Indefinido (Einführung -ar)}
\newcommand{\sequenz}{08}
\newcommand{\block}{c}
\newcommand{\fach}{Spanisch}
\newcommand{\klasse}{7}
\newcommand{\maxpunkte}{24}
\newcommand{\datum}{\today}

% ═══════════════════════════════════════════════════════════════
% KOPF- UND FUSSZEILE
% ═══════════════════════════════════════════════════════════════

\pagestyle{fancy}
\fancyhf{}
\renewcommand{\headrulewidth}{0.4pt}
\renewcommand{\footrulewidth}{0.4pt}

\lhead{\fach{} -- Klasse \klasse}
\chead{AB-\sequenz\block}
\rhead{\datum}

\lfoot{}
\cfoot{}
\rfoot{Seite \thepage\ von \pageref{LastPage}}

% ═══════════════════════════════════════════════════════════════
% BEFEHLE
% ═══════════════════════════════════════════════════════════════

\newcommand{\punktebox}[1]{\hfill\fbox{\small \_/#1}}

\newcommand{\luecke}[1]{\rule{#1}{0.4pt}}

\newtcolorbox{merkkasten}[1][]{
    colback=hellgrau,
    colframe=\fachfarbe,
    colbacktitle=white,
    coltitle=\fachfarbe,
    boxrule=1pt,
    arc=0pt,
    left=8pt, right=8pt, top=8pt, bottom=8pt,
    fonttitle=\bfseries,
    attach boxed title to top left={yshift=-2mm, xshift=5mm},
    boxed title style={boxrule=0pt, colback=white},
    title=#1
}

\newtcolorbox{vokabelkasten}{
    colback=white,
    colframe=\fachfarbe,
    boxrule=0.5pt,
    arc=0pt,
    left=5pt, right=5pt, top=5pt, bottom=5pt
}

\newtcolorbox{hinweis}{
    colback=white,
    colframe=orange,
    colbacktitle=white,
    coltitle=orange,
    boxrule=1pt,
    arc=0pt,
    left=5pt, right=5pt, top=5pt, bottom=5pt,
    fonttitle=\bfseries,
    attach boxed title to top left={yshift=-2mm, xshift=5mm},
    boxed title style={boxrule=0pt, colback=white},
    title=Tipp
}

\newenvironment{aufgabe}[1]{%
    \needspace{5\baselineskip}%
    \nopagebreak[4]%
    \vspace{10pt}
    \noindent\textbf{\textcolor{\fachfarbe}{Aufgabe #1}}
    \nopagebreak[4]%
    \vspace{5pt}
    \nopagebreak[4]%
    \begin{list}{}{\leftmargin=0pt}
    \item[]
}{%
    \end{list}
}

\newlist{teilaufgaben}{enumerate}{2}
\setlist[teilaufgaben,1]{label=\alph*), leftmargin=2em}
\setlist[teilaufgaben,2]{label=\roman*), leftmargin=2em}

\newcommand{\es}[1]{\textcolor{\fachfarbe}{\textbf{#1}}}

% ═══════════════════════════════════════════════════════════════
% DOKUMENT
% ═══════════════════════════════════════════════════════════════

\begin{document}

% ═══════════════════════════════════════════════════════════════
% KOPFBEREICH
% ═══════════════════════════════════════════════════════════════

\begin{center}
    {\Large\bfseries\textcolor{\fachfarbe}{Arbeitsblatt: \thema}}
\end{center}

\vspace{5pt}

\noindent
\begin{tabularx}{\textwidth}{|X|X|X|}
\hline
\textbf{Name:} \luecke{4cm} &
\textbf{Klasse:} \klasse &
\textbf{Datum:} \luecke{2cm} \\
\hline
\end{tabularx}

\vspace{15pt}

% ═══════════════════════════════════════════════════════════════
% MERKKASTEN: Konjugation
% ═══════════════════════════════════════════════════════════════

\begin{merkkasten}[Pretérito Indefinido: Konjugation -ar Verben]
\textbf{Das Pretérito Indefinido} beschreibt \textbf{abgeschlossene} Handlungen in der Vergangenheit.

\vspace{0.5em}

\begin{center}
\begin{tabular}{|l|c|l|}
\hline
\textbf{Person} & \textbf{Endung} & \textbf{hablar (sprechen)} \\
\hline
yo & -é & habl\luecke{1.5cm} \\
tú & -aste & habl\luecke{1.5cm} \\
él/ella/usted & -ó & habl\luecke{1.5cm} \\
nosotros/as & -amos & habl\luecke{1.5cm} \\
vosotros/as & -asteis & habl\luecke{1.5cm} \\
ellos/ellas/ustedes & -aron & habl\luecke{1.5cm} \\
\hline
\end{tabular}
\end{center}

\vspace{0.5em}

\textit{Übertrage die konjugierten Formen aus dem Lernpfad!}
\end{merkkasten}

\vspace{10pt}

% ═══════════════════════════════════════════════════════════════
% MERKKASTEN: Signalwörter
% ═══════════════════════════════════════════════════════════════

\begin{merkkasten}[Signalwörter (palabras clave)]
Diese Wörter zeigen: Es geht um die Vergangenheit!

\vspace{0.5em}

\begin{tabular}{ll}
\es{ayer} & = \luecke{3cm} \\[0.8em]
\es{el año pasado} & = \luecke{3cm} \\[0.8em]
\es{el verano pasado} & = \luecke{3cm} \\[0.8em]
\es{la semana pasada} & = \luecke{3cm} \\
\end{tabular}
\end{merkkasten}

\vspace{15pt}

% ═══════════════════════════════════════════════════════════════
% AUFGABE 1: Signalwörter zuordnen
% ═══════════════════════════════════════════════════════════════

\begin{aufgabe}{1}
Verbinde die spanischen Signalwörter mit der deutschen Bedeutung. \punktebox{4}
\end{aufgabe}

\begin{center}
\begin{tabular}{rcl}
\es{ayer} & \luecke{3cm} & letzte Woche \\[0.8em]
\es{el año pasado} & \luecke{3cm} & gestern \\[0.8em]
\es{la semana pasada} & \luecke{3cm} & letzten Sommer \\[0.8em]
\es{el verano pasado} & \luecke{3cm} & letztes Jahr \\
\end{tabular}
\end{center}

\vspace{15pt}

% ═══════════════════════════════════════════════════════════════
% AUFGABE 2: Verbformen bilden
% ═══════════════════════════════════════════════════════════════

\begin{aufgabe}{2}
Konjugiere die Verben im Pretérito Indefinido. \punktebox{6}
\end{aufgabe}

\begin{vokabelkasten}
\textbf{Verben:} hablar (sprechen) | viajar (reisen) | nadar (schwimmen) | visitar (besuchen)
\end{vokabelkasten}

\vspace{0.5em}

\noindent
a) yo / hablar: \luecke{3cm} \\[0.8em]
b) tú / viajar: \luecke{3cm} \\[0.8em]
c) él / nadar: \luecke{3cm} \\[0.8em]
d) nosotros / visitar: \luecke{3cm} \\[0.8em]
e) vosotros / hablar: \luecke{3cm} \\[0.8em]
f) ellos / viajar: \luecke{3cm} \\

\vspace{15pt}

% ═══════════════════════════════════════════════════════════════
% AUFGABE 3: Lückentext
% ═══════════════════════════════════════════════════════════════

\begin{aufgabe}{3}
Ergänze die Lücken mit der richtigen Form des Pretérito Indefinido. \punktebox{6}
\end{aufgabe}

\begin{hinweis}
Achte auf die Akzente: -é und -ó haben einen Akzent!
\end{hinweis}

\vspace{0.5em}

\noindent
\textbf{Mis vacaciones en Mallorca}

\vspace{0.5em}

\noindent
El verano pasado mi familia y yo \luecke{3cm} (viajar) a Mallorca.

\vspace{0.5em}

\noindent
Yo \luecke{3cm} (nadar) mucho en el mar.

\vspace{0.5em}

\noindent
Mi hermana \luecke{3cm} (tomar) el sol en la playa.

\vspace{0.5em}

\noindent
Nosotros \luecke{3cm} (visitar) la catedral de Palma.

\vspace{0.5em}

\noindent
Mis padres \luecke{3cm} (hablar) español con los camareros.

\vspace{0.5em}

\noindent
¿Y tú? ¿\luecke{3cm} (viajar) el verano pasado?

\vspace{15pt}

% ═══════════════════════════════════════════════════════════════
% AUFGABE 4: Eigener Urlaubsbericht
% ═══════════════════════════════════════════════════════════════

\begin{aufgabe}{4}
Schreibe mindestens 4 Sätze über deinen letzten Urlaub (oder erfinde einen Urlaub). \punktebox{8}
\end{aufgabe}

\begin{hinweis}
Verwende: ein Signalwort + verschiedene Personen + verschiedene Verben!
\end{hinweis}

\vspace{0.5em}

\begin{vokabelkasten}
\textbf{Verben:} hablar -- viajar -- nadar -- tomar el sol -- visitar

\textbf{Beispiel:} \es{El año pasado yo viajé a Italia.}
\end{vokabelkasten}

\vspace{0.5em}

\noindent
\rule{\textwidth}{0.4pt}\vspace{8pt}
\rule{\textwidth}{0.4pt}\vspace{8pt}
\rule{\textwidth}{0.4pt}\vspace{8pt}
\rule{\textwidth}{0.4pt}\vspace{8pt}
\rule{\textwidth}{0.4pt}\vspace{8pt}
\rule{\textwidth}{0.4pt}\vspace{8pt}

% ═══════════════════════════════════════════════════════════════
% PUNKTEÜBERSICHT
% ═══════════════════════════════════════════════════════════════

\vspace{20pt}
\hrule
\vspace{10pt}

\begin{flushright}
\textbf{Punkte:} \luecke{1cm} / \maxpunkte
\end{flushright}

\end{document}
